\documentclass[11pt]{article}
\input{../../shared/project_style.tex}
\rhead{Basics 48 Textbook Draft}

\begin{document}
\DocTitle{Basics 48: Energetic-Particle Distributions}{VI - Kinetic theory and energetic particles}{Fusion-alpha birth, collisional slowing down, neutral-beam anisotropy, radial gradients, normalization, and numerical sampling}

\begin{overviewbox}
\textbf{Textbook draft, version 1.0.} Alfv\'en-eigenmode stability is controlled not only by the total fast-ion density but by its distribution in radius, energy, and pitch. This chapter constructs reduced energetic-particle distributions for fusion-born alpha particles and neutral-beam ions, derives the classical slowing-down form and its normalization, explains anisotropy and pressure moments, and defines data contracts suitable for Modules 2, 3, and 5.
\end{overviewbox}

\ProjectTOC

\section{Learning objectives}
After completing this note, the reader should be able to:
\begin{itemize}[leftmargin=1.6em]
\item Distinguish birth, slowing-down, and loss-modified energetic-particle distributions.
\item Normalize isotropic distributions in speed or energy coordinates.
\item Derive the classical $1/(v^3+v_c^3)$ slowing-down shape.
\item Represent pitch anisotropy and radial gradients.
\item Compute density and pressure moments of a reduced distribution.
\item Design importance-sampling distributions for orbit and resonance calculations.
\end{itemize}

\section{Prerequisites and reading strategy}
Recommended prerequisites are Basics 1, 19, 39, 41, 42, and 44. Every formula must be paired with its integration measure; otherwise normalization is ambiguous.

\section{Coordinates for energetic particles}
Useful variables include speed $v$, kinetic energy
\begin{equation}
E=\frac{mv^2}{2},
\end{equation}
pitch cosine
\begin{equation}
\xi=\frac{v_\parallel}{v}\in[-1,1],
\end{equation}
and a magnetic-moment pitch coordinate
\begin{equation}
\lambda=\frac{\mu B_{\mathrm{ref}}}{E}.
\end{equation}
For a gyrotropic distribution,
\begin{equation}
d^3v=2\pi v^2\,dv\,d\xi.
\end{equation}
Since $dv=dE/(mv)$,
\begin{equation}
d^3v=\frac{2\pi v}{m}\,dE\,d\xi.
\end{equation}
A distribution density per unit $E$ and $\xi$ therefore differs from a Cartesian velocity-space $f$ by this Jacobian.

\section{Fusion-alpha birth source}
The D--T reaction produces a 3.5 MeV alpha particle and a 14.1 MeV neutron. A reduced local alpha source is
\begin{equation}
S_\alpha(\psi,E,\xi)=
n_D(\psi)n_T(\psi)\langle\sigma v\rangle_{DT}
\,\delta(E-E_\alpha)\,g_\xi(\xi),
\end{equation}
For the distinct D and T reactants no factor of $1/2$ is used; that combinatorial factor is required only for reactions between identical species. For isotropic birth,
\begin{equation}
g_\xi(\xi)=\frac{1}{2},\qquad -1\le\xi\le1.
\end{equation}
The spatial source is strongly peaked where density and temperature make the fusion reactivity largest.

\section{Classical slowing-down balance}
At high speed, alpha particles lose energy primarily to electrons; at lower speed, ion drag becomes important. A simplified isotropic steady kinetic balance in speed is
\begin{equation}
\frac{1}{v^2}\frac{d}{dv}\left[v^2\dot v\,f(v)\right]
+S_0\frac{\delta(v-v_\alpha)}{4\pi v^2}=0.
\end{equation}
A standard drag model has
\begin{equation}
\dot v=-\frac{\nu_s}{v^2}(v^3+v_c^3),
\end{equation}
where $v_c$ is the critical speed separating electron-dominated and ion-dominated slowing. Below the birth speed and away from the delta source, the speed-space flux is constant, giving
\begin{equation}
\boxed{f_s(v)=\frac{C}{v^3+v_c^3}H(v_\alpha-v)}.
\end{equation}
This is a non-Maxwellian distribution with a broad energetic tail.

\section{Normalization of the slowing-down distribution}
For an isotropic population of density $n_f$,
\begin{equation}
n_f=4\pi\int_0^{v_\alpha} f_s(v)v^2\,dv.
\end{equation}
Using
\begin{equation}
\int\frac{v^2}{v^3+v_c^3}\,dv
=\frac{1}{3}\ln(v^3+v_c^3),
\end{equation}
we obtain
\begin{equation}
n_f=\frac{4\pi C}{3}
\ln\!\left[1+\left(\frac{v_\alpha}{v_c}\right)^3\right].
\end{equation}
Therefore
\begin{equation}
\boxed{C=\frac{3n_f}{4\pi
\ln\!\left[1+(v_\alpha/v_c)^3\right]}}.
\end{equation}
This normalization is valid for the six-dimensional velocity distribution $f_s(v)$ integrated with $4\pi v^2dv$.

\section{Energy-space form}
Define an energy density $F_E(E)$ such that
\begin{equation}
n_f=\int_0^{E_\alpha}F_E(E)\,dE.
\end{equation}
For an isotropic velocity distribution,
\begin{equation}
F_E(E)=4\pi v^2f(v)\frac{dv}{dE}
=\frac{4\pi v}{m}f(v).
\end{equation}
Confusing $f(v)$ with $F_E(E)$ changes both units and shape. Machine-readable outputs should store coordinate names and measures explicitly.

\section{Neutral-beam distributions}
Neutral-beam injection produces a directional source with finite energy spread. A simple model is
\begin{equation}
S_{\mathrm{NBI}}(\psi,E,\xi)=S_0g_r(\psi)
\exp\!\left[-\frac{(E-E_b)^2}{2\Delta E^2}\right]
\exp\!\left[-\frac{(\xi-\xi_b)^2}{2\Delta\xi^2}\right].
\end{equation}
Co-current tangential injection has $\xi_b$ near one, while perpendicular injection has smaller $|\xi_b|$. Ionization geometry and orbit topology make the true source nonseparable; the Gaussian model is only a controlled benchmark.

\section{Radial profiles and gradients}
A reduced separable distribution may be written
\begin{equation}
f_f(\psi,E,\xi)=n_f(\psi)\,\widehat f_E(E;\psi)
\,\widehat f_\xi(\xi;\psi),
\end{equation}
where each hatted factor is normalized under its stated measure. Example radial profiles include
\begin{equation}
n_f(s)=n_{f0}(1-s^a)^b,
\qquad 0\le s\le1.
\end{equation}
The logarithmic gradient
\begin{equation}
\frac{d\ln n_f}{ds}
\end{equation}
is often more relevant to wave drive than the absolute density. Alfv\'en-eigenmode transport tends to flatten resonant gradients.

\section{Pressure and anisotropy moments}
For a gyrotropic distribution,
\begin{align}
p_\parallel&=m\int v_\parallel^2 f\,d^3v,\\
p_\perp&=\frac{m}{2}\int v_\perp^2 f\,d^3v.
\end{align}
An isotropic population has $p_\parallel=p_\perp=p_f$. A beam-like distribution can have strong anisotropy, changing both equilibrium pressure and kinetic coupling to modes. The fast-ion beta is
\begin{equation}
\beta_f=\frac{2\mu_0p_f}{B^2}
\end{equation}
for isotropic scalar pressure.

\section{Loss cones and orbit filtering}
The source distribution is not necessarily the confined distribution. Prompt orbit losses can remove selected pitch and radial regions. A reduced loss-cone factor might be
\begin{equation}
L(\psi,E,\xi)\in[0,1],
\qquad
f_{\mathrm{conf}}=L f_{\mathrm{source}}.
\end{equation}
Module 3 should estimate $L$ from orbit following rather than assuming it. This creates the coupling
\begin{equation}
S_f\longrightarrow\text{orbits}\longrightarrow f_{\mathrm{conf}}
\longrightarrow\text{mode drive}.
\end{equation}

\section{Sampling and marker weights}
If markers are sampled from $g(\psi,E,\xi)$, an integral is estimated by weights $w=f/g$. Uniform sampling is inefficient when resonance occupies a narrow phase-space region. Useful strategies include:
\begin{itemize}[leftmargin=1.6em]
\item sample proportional to particle density for global moments;
\item oversample high-energy tails for wall-loss power;
\item oversample near resonance surfaces for kinetic drive;
\item use stratified sampling in pitch and radius to capture rare trapped or loss orbits.
\end{itemize}
Weights and marker-density normalization must be stored with the particle data.

\section{Uncertainty and model hierarchy}
A parameterized slowing-down distribution is appropriate for early verification, but predictive studies require collisions, sources, orbit loss, radial transport, and possibly time dependence:
\begin{equation}
\frac{\partial f_f}{\partial t}
+\dot{\mathbf Z}\cdot\nabla_{\mathbf Z}f_f
=C[f_f]+S_f-L_f.
\end{equation}
The hierarchy should separate uncertainty in source rate, critical speed, radial profile, pitch dependence, and loss model.

\section{Connection to the Kinetic MHD-PINN project}
Module 2 owns the source and distribution representation. Module 3 applies orbit filtering and computes frequencies. Module 5 evaluates resonant drive using gradients of $f_f$. Module 6 evolves redistribution. The shared-state contract should include coordinate definitions, Jacobian, units, normalization, species, reference magnetic field, and time stamp.

\section{Computational laboratory}
\begin{enumerate}[leftmargin=1.6em]
\item Normalize the classical slowing-down distribution by quadrature and compare with the analytic constant $C$.
\item Transform it from speed to energy coordinates and verify that both integrals yield $n_f$.
\item Sample markers by inverse or rejection sampling and recover density and pressure moments.
\item Add a pitch Gaussian and measure $p_\parallel/p_\perp$.
\item Apply a synthetic loss cone and quantify changes in density, energy, and resonant population.
\end{enumerate}

\section{Summary}
Energetic-particle stability depends on the full distribution in space, energy, and pitch. Fusion alphas are born nearly monoenergetically and relax toward a classical slowing-down form; neutral beams produce anisotropic sources. Correct Jacobians and normalization are indispensable. Radial and phase-space gradients supply free energy for Alfv\'en eigenmodes, while orbit loss and transport reshape those gradients.

\section{Exercises}
\begin{enumerate}[leftmargin=1.6em]
\item Derive the normalization constant for $f_s(v)=C/(v^3+v_c^3)$.
\item Derive the energy-coordinate density $F_E(E)$.
\item Show that an isotropic distribution satisfies $p_\parallel=p_\perp$.
\item Explain why total fast-ion density alone cannot determine AE stability.
\item Design a JSON schema that prevents ambiguity between $f(v)$ and $F_E(E)$.
\end{enumerate}

\SymbolsAcronymsSection
\begin{symtable}
$f_f$ & Energetic-particle distribution & Fast-ion phase-space density under a stated measure.\\
$E_\alpha$ & Alpha birth energy & Approximately $3.5$ MeV for D--T fusion.\\
$v_c$ & Critical speed & Speed separating electron- and ion-dominated collisional slowing.\\
$\xi$ & Pitch cosine & $v_\parallel/v$.\\
$\lambda$ & Magnetic pitch coordinate & Normalized magnetic moment coordinate.\\
$\beta_f$ & Fast-ion beta & Ratio of fast-ion pressure to magnetic pressure.\\
\end{symtable}

\section{References}
\begin{enumerate}[leftmargin=1.6em]
\item J. D. G. Cordey and M. J. Houghton, classical fast-ion slowing-down literature.
\item R. B. White, \emph{The Theory of Toroidally Confined Plasmas}.
\item J. Wesson, \emph{Tokamaks}.
\item L. Chen and F. Zonca, energetic-particle and Alfv\'en-wave review literature.
\end{enumerate}

\end{document}
